%! Sine, Cosine, and Tangent — Unit 6, Lesson 6.3 — Guided Notes
\documentclass[10pt]{article}
\usepackage{precalculus-article}
\usepackage{precalculus-boxes}
\newcommand{\TallMath}[1]{$\displaystyle #1\rule[-1.4em]{0pt}{3.2em}$}

\begin{document}

\pageheader{Unit 6, Lesson 6.3}{Guided Notes}
\namedateperiod

\vspace{0.10in}

%--- Objective ---
\begin{objectivebox}
\textbf{LO 3.2.A:} Determine the sine, cosine, and tangent of an angle using the unit circle.\\[4pt]
\writelines{2}
\end{objectivebox}

\vspace{0.10in}

%--- Vocabulary ---
\begin{vocabbox}
\termblanklong{standard position}
\termblanklong{terminal ray}
\termblanklong{coterminal angles}
\termblanklong{radian measure}
\termblanklong{unit circle}
\termblanklong{sine ($\sin\theta$)}
\termblanklong{cosine ($\cos\theta$)}
\termblanklong{tangent ($\tan\theta$)}
\end{vocabbox}

\vspace{0.10in}

%--- Hook ---
\begin{hookbox}
\textbf{Spinning Wheel:} A point $P$ sits on the rim of a wheel of radius 1, starting at $(1,\, 0)$.
After the wheel rotates counterclockwise by angle $\theta$, the point is at position $(\underline{\hspace{1.2cm}},\, \underline{\hspace{1.2cm}})$.
\begin{itemize}[leftmargin=1.4em, itemsep=2pt]
  \item How far right/left of center is $P$? \hfill The answer is called: \blank{2.4cm}
  \item How far up/down from center is $P$? \hfill The answer is called: \blank{2.4cm}
  \item What is the slope of the line from center to $P$? \hfill The answer is called: \blank{2.4cm}
\end{itemize}
\end{hookbox}

\vspace{0.10in}

%--- Radians ---
\begin{notesbox}{Radian Measure}
\textbf{Definition:} Radian measure $=$ \blank{3cm} $\div$ \blank{3cm}

\medskip
On a \textbf{unit circle} ($r = 1$): radian measure $=$ \blank{3.5cm}

\medskip
\begin{center}
\renewcommand{\arraystretch}{1.6}
\begin{tabular}{|c|c|c|}
\hline
\textbf{Degrees} & \textbf{Radians} & \textbf{Arc length (unit circle)} \\
\hline
$360°$  & \blank{1.8cm} & \blank{1.8cm} \\
$180°$  & \blank{1.8cm} & \blank{1.8cm} \\
$90°$   & \blank{1.8cm} & \blank{1.8cm} \\
\hline
\end{tabular}
\end{center}

\medskip
Conversion: $\pi$ radians $=$ \blank{1.8cm} degrees \qquad $\dfrac{\pi}{2}$ radians $=$ \blank{1.8cm} degrees
\end{notesbox}

\vspace{0.10in}

%--- Unit Circle Definitions ---
\begin{notesbox}{LO 3.2.A --- The Unit Circle}
Let $\theta$ be an angle in standard position. Let $P = (x,\, y)$ be the point where the
\textbf{terminal ray} meets the \textbf{unit circle}.

\medskip
\begin{center}
\renewcommand{\arraystretch}{2.0}
\begin{tabular}{|c|c|c|}
\hline
\textbf{Function} & \textbf{Unit-Circle Value} & \textbf{General Definition (radius $r$)} \\
\hline
$\sin\theta$ & \blank{2.8cm} & \blank{2.8cm} \\
$\cos\theta$ & \blank{2.8cm} & \blank{2.8cm} \\
$\tan\theta$ & \blank{2.8cm} & \blank{2.8cm} \\
\hline
\end{tabular}
\end{center}

\medskip
\textbf{Pythagorean Identity:} Since $x^2 + y^2 = 1$ on the unit circle:
\[
  \underline{\hspace{1.8cm}}^2 + \underline{\hspace{1.8cm}}^2 = 1
\]
\end{notesbox}

\vspace{0.10in}

%--- Practice ---
\begin{practicebox}
\textbf{Practice: Evaluate using the unit circle.} For each angle, identify the point $P = (x,y)$
on the unit circle, then fill in the table.

\medskip
\begin{center}
\renewcommand{\arraystretch}{2.0}
\begin{tabular}{|c|c|c|c|c|}
\hline
$\theta$ & Point $P = (x,y)$ & $\cos\theta$ & $\sin\theta$ & $\tan\theta$ \\
\hline
$0$           & \blank{2.8cm} & \blank{1.6cm} & \blank{1.6cm} & \blank{1.6cm} \\
$\pi/2$       & \blank{2.8cm} & \blank{1.6cm} & \blank{1.6cm} & \blank{1.6cm} \\
$\pi$         & \blank{2.8cm} & \blank{1.6cm} & \blank{1.6cm} & \blank{1.6cm} \\
$3\pi/2$      & \blank{2.8cm} & \blank{1.6cm} & \blank{1.6cm} & \blank{1.6cm} \\
$2\pi$        & \blank{2.8cm} & \blank{1.6cm} & \blank{1.6cm} & \blank{1.6cm} \\
\hline
\end{tabular}
\end{center}
\end{practicebox}

\end{document}
